\documentclass{abnt}
\usepackage[utf8]{inputenc}
\usepackage[T1]{fontenc}
\usepackage{geometry}
\usepackage{graphicx}
\usepackage{float}
\usepackage{hyperref}
\usepackage[english,indonesian]{babel}

\title{Optimasi Pemrosesan Template LaTeX Menggunakan Pendekatan Berbasis AI}
\author{Academic Support Agent}
\date{\today}

\begin{document}
\maketitle

\begin{abstract}

\end{abstract}

\chapter{PENDAHULUAN}


\chapter{TINJAUAN PUSTAKA}


% Original content (commented out):
% % Pre-ambulo
\documentclass[a4paper, 12pt]{abnt}

\usepackage[english]{babel}
\usepackage[utf8]{inputenc}
\usepackage[T1]{fontenc}
\usepackage{dsfont}
\usepackage{amssymb,amsmath}
\usepackage{multirow}
\usepackage[alf]{abntcite}
\usepackage[pdftex]{color, graphicx}
\usepackage{colortbl}
\usepackage{url}
\usepackage{abnt-alf}
\usepackage{abntcite}
\usepackage{algorithm}
\usepackage{algorithmic}

% Pacotes adicionados
\usepackage{amsthm}
\usepackage{booktabs}
\usepackage{caption}
\usepackage{subcaption}
\usepackage{longtable}
\usepackage{array}
\usepackage{tikz}
\usetikzlibrary{shapes,arrows}
\usepackage{filecontents,pgfplots}
\usepackage{pgfplotstable}
\usepackage[inline]{enumitem}
\usetikzlibrary{positioning, arrows}

\usepackage[protrusion=true,expansion=true]{microtype}

\theoremstyle{definition}
\newtheorem{definition}{Definition}[section]
 
\theoremstyle{remark}
\newtheorem*{remark}{Remark}
\usepackage{arydshln}
%\usepackage{alg}
%\usepackage{hyperref}


% Redefinicao de instrucoes
\floatname{algorithm}{Algoritmo}
\renewcommand{\algorithmicrequire}{\textbf{Entrada:}}
\renewcommand{\algorithmicensure}{\textbf{Saída:}}
\renewcommand{\algorithmicend}{\textbf{fim}}
\renewcommand{\algorithmicif}{\textbf{se}}
\renewcommand{\algorithmicthen}{\textbf{então}}
\renewcommand{\algorithmicelse}{\textbf{senão}}
\renewcommand{\algorithmicfor}{\textbf{para}}
\renewcommand{\algorithmicforall}{\textbf{para todo}}
\renewcommand{\algorithmicdo}{\textbf{faça}}
\renewcommand{\algorithmicwhile}{\textbf{enquanto}}
\renewcommand{\algorithmicloop}{\textbf{loop}}
\renewcommand{\algorithmicrepeat}{\textbf{repetir}}
\renewcommand{\algorithmicuntil}{\textbf{até  que}}
\renewcommand{\algorithmiccomment}[1]{\% #1}


% Definicao da lista de simbolos
% \simb[entrada na lista de simbolos]{simbolo}:
% Escreve o simbolo no texto e uma entrada na lista de simbolos.
% Se o parametro opcional e omitido, usa-se o parametro obrigatorio.
\newcommand{\simb}[2][]
{%
	\ifthenelse{\equal{#1}{}}
	{\addcontentsline{los}{simbolo}{#2}}
	{\addcontentsline{los}{simbolo}{#1}}
}
% Para aceitar comandos com @ (at) no nome
\makeatletter 
% \listadesimbolos: comando que imprime a lista de simbolos
\newcommand{\listofsymbols}
{
	\pretextualchapter{List of symbols}
	{\setlength{\parindent}{0cm}
	\@starttoc{los}}
}
% Como a entrada sera impressa
\newcommand\l@simbolo[2]{\par #1}
\makeatother


% Definicao da lista de abreviaturas e siglas
% \abrv[entrada na lista de simbolos]{abreviatura}:
% Escreve a sigla/abreviatura no texto e uma entrada na lista de abreviaturas e siglas.
% Se o parametro opcional e omitido, usa-se o parametro obrigatorio.
\newcommand{\abrv}[2]
{
	\ifthenelse{\equal{#1}{}}
	{\addcontentsline{loab}{abreviatura}{#2}}
	{\addcontentsline{loab}{abreviatura}{#1}}#2%
}

% Para aceitar comandos com @ (at) no nome
\makeatletter 
% \listadeabreviaturas: comando que imprime a lista de abreviaturas e siglas
\newcommand{\listadeabreviaturas}
{
	\pretextualchapter{List of abbreviations}
	{\setlength{\parindent}{0cm}
	\@starttoc{loab}}
}
% Como a entrada sera impressa
\newcommand\l@abreviatura[2]{\par #1}
%\newcommand\l@abreviatura[2]{} % Não imprime a entrada, apenas adiciona na tabela
\makeatother


% \listofalgorithms: comando que imprime a lista de algoritmos
\renewcommand{\listalgorithmname}{List of algorithms}


% Hifeniza  o de palavras feita de forma incorreta pelo LaTeX
\hyphenation{PYTHON ou-tros}


% Inicio do documento
\begin{document}

	\frenchspacing
	
	% Capa (arquivo Includes/Capa.tex)
	\include{Includes/Capa}

	% Folha de rosto (arquivo Includes/FolhaRosto.tex)
	\include{Includes/FolhaRosto}	
	
	% Ficha catalográfica
	\include{Includes/FichaCatalografica}
	
	% Folha de aprovacao (arquivo Includes/FolhaAprovacao.tex)
	\include{Includes/FolhaAprovacao}	
	
	% Dedicatoria (arquivo Includes/Dedicatoria.tex)
	\include{Includes/Dedicatoria}
	
	% Agradecimentos (arquivo Includes/Agradecimentos.tex)
	\include{Includes/Agradecimentos}
   
   % Epigrafe (arquivo Includes/Epigrafe.tex)
	\include{Includes/Epigrafe}
	
	% Resumo em l ngua vernacula (arquivo Includes/Resumo.tex)
%	\include{Includes/Resumo}
	
	% Abstract, resumo em l ngua estrangeira (arquivo Include/Abstract.tex)
	\include{Includes/Abstract}

	% Sum rio
	\tableofcontents
	
	% Lista de figuras
	\listoffigures

	% Lista de tabelas
	\listoftables
	
	% Lista de abreviaturas e siglas
	\listadeabreviaturas
	
	% Lista de s mbolos
%	\listofsymbols
	
	% Lista de algoritmos (se houver)
	% Devem ser inclu dos os pacotes algorithm e algorithmic
	% \listofalgorithms
	


	% Parte central do trabalho, englobando os cap tulos que constituem o mesmo
	% Os referidos cap tulos devem ser organizados dentro do diret rio "Cap tulos"

	% Capitulo 1: Introdu  o (arquivo Includes/Introducao.tex)
	\include{Capitulos/Introducao}
	
	% Capitulo 2: Segundo cap tulo (arquivo Includes/Capitulo2.tex)
	\include{Capitulos/Capitulo2}
	
	% Capitulo 3: Terceiro cap tulo (arquivo Includes/Capitulo3.tex)
	\include{Capitulos/Capitulo3}
	
	% Capitulo 4: Quarto cap tulo (arquivo Includes/Capitulo4.tex)
	\include{Capitulos/Capitulo4}
	
	% Capitulo 5: Quinto cap tulo (arquivo Includes/Capitulo5.tex)
	\include{Capitulos/Capitulo5}
		
	% Consideracoes finais
	\include{Capitulos/Consideracoes}
	
	% Bibliografia (arquivo Capitulos/Referencias.bib)
	\bibliography{Capitulos/Referencias}
	\bibliographystyle{abnt-alf}
	
	% Ap ndice A (arquivo Includes/ApendiceA)
%	\include{Capitulos/ApendiceA}
	
	% Anexo A (arquivo Includes/AnexoA)
%	\include{Capitulos/AnexoA}
	
	% P gina em branco
	%\newpage
	%\include{Includes/AbreviturasOrdenadas}

\end{document}


\end{document}